\documentclass[10pt,letterpaper]{article}
\usepackage{cornell-notes}

\title{Clause 29.07.08: Class High Resolution Clock}
\author{}
\date{}
\setDocTitle{Clause 29.07.08: Class High Resolution Clock}
\setDocSubtitle{C++ 2024}
\setDocOwner{C++ 2024 Cornell Notes}
\setDocDate{\today}
\setCornellCollection{C++ 2024 Cornell Notes}
\setCornellUnitType{Clause}
\setCornellUnitNumber{29.07.08}
\setCornellUnitTitle{Class High Resolution Clock}
\hypersetup{pdftitle={Clause 29.07.08: Class High Resolution Clock},pdfsubject={Cornell notes},pdfauthor={}}

\begin{document}
\maketitle
\begin{CornellOverviewBox}{Overview}
\textbf{Classification:} Standard library - Normative\\[2pt]
\textbf{Learning objective:} Understand the rules, terminology, and semantic consequences associated with Class high\_resolution\_clock.
\end{CornellOverviewBox}\begin{CornellSummaryBox}{Summary}
{\sffamily\bfseries\color{Accent} Summary}\par\vspace{3pt}Section 29.7.8 focuses on Class high\_resolution\_clock. Use the listed grammar productions together with the semantic constraints and disambiguation rules referenced by the section. When applying it, preserve the distinction between mandatory requirements, recommendations, permissions, and behavior deliberately left undefined, unspecified, or implementation-defined.
\end{CornellSummaryBox}

\end{document}
